\section{Perspectivas futuras y la frontera de la IA}

\subsection{¿RL for LLM todavía necesita un Breakthrough?}

La respuesta de Weng Jiayi (翁家翌) es que \textbf{es posible}, pero la tarea más urgente es primero exprimir al máximo los métodos y los recursos de cómputo existentes. El estado actual es que «todavía no se ha hecho scale up por completo»: primero hay que hacer hill climb, ver cuál es el límite superior de los métodos existentes, y después considerar si se necesita un nuevo paradigma.

Él enfatizó una idea clave: \textbf{no se puede usar el estado actual para predecir qué sucederá después}. Podría haber un nuevo paradigma de RL, podría haber un nuevo paradigma de post-training, cualquiera de los dos es posible. Cada día se enfrenta a retos desconocidos.

\subsection{¿Dónde está el mayor cuello de botella?}

Weng Jiayi considera que el mayor cuello de botella de los modelos grandes en el futuro sigue siendo la infra:

\begin{enumerate}
    \item El \textbf{rendimiento} para reparar la infra: cuántos bugs se pueden corregir por unidad de tiempo
    \item Cuántas veces se puede \textbf{iterar correctamente} por unidad de tiempo
    \item Estos dos indicadores determinan la eficiencia productiva de todo el equipo y pueden potenciar todo el resto del trabajo
\end{enumerate}

«Ya sea el algoritmo o el entorno: si arreglas todos los bugs, es posible que el algoritmo funcione muy bien sin necesidad de cambiar nada.»

\subsection{Si la IA pudiera resolver un gran problema del mundo}

Cuando le preguntaron qué gran problema del mundo le gustaría que la IA resolviera, la respuesta de Weng Jiayi fue sorprendente: \textbf{cómo predecir el futuro}, no predecir cómo cae una taza, sino predecir toda una vida, la configuración del mundo, todo. Esta respuesta es coherente con su visión determinista del mundo: si el mundo es un proceso de Márkov determinista, en teoría debería poder predecirse.

Pero al mismo tiempo también es consciente de que una herramienta así podría ser un desastre: «Si se consigue una máquina capaz de predecir el futuro, eso sería en realidad un desastre para cada persona: provocaría el colapso de todo el sistema de valores.» Incluso considera que, si llegara a desarrollarse un modelo de IA así, «la mejor opción sería destruirlo, para que nunca llegara a salir».

\subsection{La rutina diaria y el cuidado del cuerpo de Weng Jiayi}

Weng Jiayi reveló cómo es su estado de trabajo en OpenAI. Considera que su trabajo de mantenimiento cotidiano «en realidad no requiere tanto coeficiente intelectual»: basta con hacer lo correcto en la dirección correcta. También compartió un detalle muy contrastante: en Tsinghua (清华) reprobó los 3,000 metros en la clase de educación física, pero ahora tiene el hábito de correr 3,000 metros dos veces por semana: «Antes que nada, hay que asegurarse de que el cuerpo esté sano.» Esto es una extensión de su filosofía de «invertir en el futuro» aplicada al cuidado del cuerpo.

\begin{warningbox}{La humildad y la autoconciencia de un genio}
Weng Jiayi enfatizó varias veces que él no es especial: «Si me sustituyeran por cualquier otra persona, si esa persona tuviera mi context, debería poder desempeñarse perfectamente bien también». Considera que tiene mucha suerte de estar en esta posición, pero que esto «lo podría hacer cualquier ser humano normal». Esta humildad podría ser sincera: él atribuye su éxito más a \textbf{el lugar correcto} (OpenAI + RL Infra) y a \textbf{un context suficiente}, y no al talento personal.
\end{warningbox}

\subsection{Resumen de la sección}

Para Weng Jiayi, el futuro de la IA está lleno tanto de certidumbre (la ruta del scale up es clara, la infra necesita optimizarse constantemente) como de incertidumbre (un nuevo paradigma podría surgir en cualquier momento). La estrategia que elige para enfrentarlo es la de siempre: hacer lo correcto en la dirección correcta, dejar la infra bien construida y permitir que la iteración rápida lo resuelva todo.

